\chapter{Funciones}\label{ch:g10:functions}

A \omterm{def:g10:functions:function}{función} es una máquina que toma un número y produce un número,
Siempre la misma salida para la misma entrada. funciones son los centrales
Objetos del análisis que construirás en los próximos años.
(\cref{ch:g11:quad} adelante); este capítulo establece el vocabulario ---
\omterm{def:g10:functions:function}{dominio}, \omterm{def:g10:functions:function}{imagen}, grafo, variaciones --- y entrena la habilidad esencial de
leer información de un grafo.

\section{Vocabulario: imágenes y preimágenes}

\begin{definition}[Función]\label{def:g10:functions:function}
Sea $D $ ser un conjunto de números reales. A \emph{función}\index{función} $ f $
definido en $ D $ asociados a cada número $ x \in D $ exactamente uno real
numero, escrito $ f(x)$ y llamó al \emph{imagen}\index{imagen} de $ x $.
el conjunto $ D $ es el \emph{dominio}\index{dominio} de $ f $. escribimos
\[
f : x \longmapsto f(x).
\]
Si $ f(x) = y $, entonces $ x $ es un \emph{preimagen}\index{preimagen} de $ y $.
\end{definition}

\begin{remark}
Cada $ x $ en el \omterm{def:g10:functions:function}{dominio} tiene exactamente uno \omterm{def:g10:functions:function}{imagen}, pero un número $ y $ puede tener
varios preimágenes, o ninguno. Para $ f(x) = x^2$: el \omterm{def:g10:functions:function}{imagen} de $3$ es $9$, y
$9$ tiene dos preimágenes, $3$ y $-3$; el numero $-4$ no tiene \omterm{def:g10:functions:function}{preimagen}.
\end{remark}

\begin{example}[Encontrar un dominio]\label{ex:g10:functions:domain}
cuando una \omterm{def:g10:functions:function}{función} viene dada por una fórmula, su \omterm{def:g10:functions:function}{dominio} es el conjunto de $ x $ para
que la fórmula tiene sentido.
\begin{itemize}
  \item $ f(x) = 3x^2 - 5x + 1$ tiene sentido para cada $ x $: el \omterm{def:g10:functions:function}{dominio} es
        $\R $.
  \item $ g(x) = \dfrac{1}{x - 2}$ requiere $ x - 2 \neq 0$: el \omterm{def:g10:functions:function}{dominio} es
        todos los reales excepto $2$.
  \item $ h(x) = \sqrt{x - 3}$ requiere $ x - 3 \geq 0$: el \omterm{def:g10:functions:function}{dominio} es
        $\intco{3}{+\infty}$.
\end{itemize}
\end{example}

\begin{example}[Computación de imágenes y preimágenes.]\label{ex:g10:functions:images}
Sea $ f(x) = x^2 - 4x + 3$, definido en $\R $.

\emph{\omterm{def:g10:functions:function}{Imagen} de $5$:} sustituto $ x = 5$:
$ f(5) = 25 - 20 + 3 = 8$.

\emph{Preimágenes de $3$:} resolver $ f(x) = 3$:
\[
x^2 - 4x + 3 = 3
\quad\Longleftrightarrow\quad
x^2 - 4x = 0
\quad\Longleftrightarrow\quad
x(x - 4) = 0,
\]
entonces el preimágenes de $3$ son $0$ y $4$.
\end{example}

\section{La gráfica de una función.}

\begin{definition}[Gráfico]\label{def:g10:functions:graph}
En un sistema de coordenadas, el \emph{gráfico}\index{gráfico} (o \emph{curva}) de
$ f $ es el conjunto de todos los puntos $(x, f(x))$ para $ x $ en el \omterm{def:g10:functions:function}{dominio}. en otros
palabras, un punto $(x, y)$ se encuentra en el gráfico exactamente cuando $ y = f(x)$.
\end{definition}

\begin{omfigure}
\begin{tikzpicture}
\begin{axis}[omaxis,
  width=8.6cm, height=6.0cm,
  xmin=-2.6, xmax=4.6, ymin=-3.4, ymax=4.4,
  xtick={-2,-1,1,2,3,4}, ytick={-3,-2,-1,1,2,3,4}]
  \addplot[omDef, thick, domain=-2.2:4.2, samples=120]
    {0.25*x^3 - 0.75*x^2 - x + 2};
  \draw[dashed, omThm] (axis cs:3,0) -- (axis cs:3,1.25)
    -- (axis cs:0,1.25);
  \fill[omThm] (axis cs:3,1.25) circle (1.5pt);
  \node[omThm, font=\small, above right] at (axis cs:3,1.2)
    {$(3,\,f(3))$};
\end{axis}
\end{tikzpicture}

{\small leyendo un \omterm{def:g10:functions:function}{imagen} en un grafo: empezar desde $ x = 3$ en la horizontal
eje, vaya verticalmente a la curva, luego horizontalmente al eje vertical
leer $ f(3) = 1.25$.}
\end{omfigure}

\begin{method}[Leyendo un gráfico]\label{met:g10:functions:reading}
en el grafo de $ f $:
\begin{enumerate}
  \item \emph{\omterm{def:g10:functions:function}{imagen} de $ a $}: ir verticalmente desde $ x = a $ en la horizontal
        eje a la curva, luego horizontalmente al eje vertical; el
        valor leído hay $ f(a)$;
  \item \emph{preimágenes de $ b $}: dibuja la línea horizontal $ y = b $; el
        preimágenes son los $ x $-coordenadas de todos sus \omterm{def:g10:numbers:interunion}{intersección} puntos
        con la curva;
  \item \emph{soluciones de $ f(x) = k $}: igual que el preimágenes de $ k $;
  \item \emph{soluciones de $ f(x) \leq k $}: el $ x $ para lo cual la curva es
        en o debajo de la línea $ y = k $.
\end{enumerate}
\end{method}

\begin{omfigure}
\begin{tikzpicture}
\begin{axis}[omaxis,
  width=8.6cm, height=6.0cm,
  xmin=-3.4, xmax=3.4, ymin=-2.4, ymax=4.4,
  xtick={-2,2}, xticklabels={$ x_1$,$ x_2$}, ytick={2}, yticklabels={$ k $}]
  \addplot[omDef, thick, domain=-3:3, samples=120] {0.5*x^2};
  \addplot[omThm, thick, domain=-3.2:3.2] {2};
  \fill (axis cs:-2,2) circle (1.5pt);
  \fill (axis cs:2,2) circle (1.5pt);
  \draw[dashed] (axis cs:-2,2) -- (axis cs:-2,0);
  \draw[dashed] (axis cs:2,2) -- (axis cs:2,0);
\end{axis}
\end{tikzpicture}

{\small resolviendo $ f(x) = k $ gráficamente: las soluciones $ x_1$ y $ x_2$ son
las abscisas de la \omterm{def:g10:numbers:interunion}{intersección} puntos de la curva con la horizontal
linea $ y = k $. Aquí $ f(x) \leq k $ aguanta $\intcc{x_1}{x_2}$, donde el
la curva está debajo de la línea.}
\end{omfigure}

\begin{example}\label{ex:g10:functions:membership}
es el punto $ A(2, 5)$ en el grafo de $ f(x) = x^2 + 1$? Calcular
$ f(2) = 4 + 1 = 5$: si, desde $ f(2)$ es igual al $ y $-coordenada de $ A $.
el punto $ B(3, 8)$ no está en el grafo, porque $ f(3) = 10 \neq 8$.
\end{example}

\section{Variaciones de una función}

\begin{definition}[Creciendo, disminuyendo]\label{def:g10:functions:variations}
Sea $f$ definida en un \omterm{def:g10:numbers:interval}{intervalo} $ I $.
\begin{itemize}
  \item $ f $ es \emph{creciente}\index{creciente} en $ I $ cuando
        para todos $ u, v \in I $, si $ u < v $ entonces $ f(u) \leq f(v)$: el
        Los productos crecen con los insumos y el grafo sube de izquierda a
        correcto.
  \item $ f $ es \emph{decreciente}\index{decreciente} en $ I $ cuando
        $ u < v $ implica $ f(u) \geq f(v)$: el grafo cae de izquierda a
        correcto.
\end{itemize}
Con desigualdades estrictas $ f(u) < f(v)$ (resp.\ $>$) decimos
\emph{estrictamente} creciente (resp.\ decreciente).
\end{definition}

\begin{definition}[tabla de variaciones]\label{def:g10:functions:table}
A \emph{tabla de variaciones}\index{tabla de variaciones} resume sobre qué
\omterm{def:g10:numbers:interval}{intervalos} $ f $ aumenta y disminuye, con flechas $\nearrow $ y
$\searrow $ y registra los valores de $f $ en los puntos de inflexión.
\end{definition}

\begin{example}\label{ex:g10:functions:table}
Considere la \omterm{def:g10:functions:function}{función} graficado a continuación en $\intcc{-2}{4}$.

\begin{omfigure}
\begin{tikzpicture}
\begin{axis}[omaxis,
  width=8.6cm, height=5.4cm,
  xmin=-2.8, xmax=4.6, ymin=-2.4, ymax=3.8,
  xtick={-2,1,4}, ytick={-1.5,3}, yticklabels={$-1.5$,$3$}]
  \addplot[omDef, thick, domain=-2:4, samples=120] {-0.5*(x-1)^2 + 3};
  \fill (axis cs:-2,-1.5) circle (1.5pt);
  \fill (axis cs:4,-1.5) circle (1.5pt);
  \fill (axis cs:1,3) circle (1.5pt);
\end{axis}
\end{tikzpicture}

{\small A \omterm{def:g10:functions:function}{función} definido en $\intcc{-2}{4}$, \omterm{def:g10:functions:variations}{creciente} entonces \omterm{def:g10:functions:variations}{decreciente}.}
\end{omfigure}

leyendo el grafo: $ f $ aumenta de $ f(-2) = -1.5$ hasta su \omterm{def:g10:functions:extrema}{máximo}
$ f(1) = 3$, luego disminuye hasta $ f(4) = -1.5$. Es \omterm{def:g10:functions:table}{tabla de variaciones} es
\begin{center}
\renewcommand{\arraystretch}{1.3}
\begin{tabular}{c|ccccc}
$ x $ & $-2$ & & $1$ & & $4$ \\
\hline
$ f $ & $-1.5$ & $\nearrow $ & $3$ & $\searrow $ & $-1.5$ \\
\end{tabular}
\end{center}
\end{example}

\begin{example}[Demostrando una variación]\label{ex:g10:functions:proof}
mostrar que $ f(x) = 3x + 1$ es estrictamente \omterm{def:g10:functions:variations}{creciente} en $\R $. toma cualquier
$ u < v $ y comparar el imágenes:
\[
f(v) - f(u) = (3v + 1) - (3u + 1) = 3(v - u) > 0,
\]
porque $ v - u > 0$. Entonces $ f(u) < f(v)$: la \omterm{def:g10:functions:function}{función} es estrictamente
\omterm{def:g10:functions:variations}{creciente}. El mismo cálculo con pendiente negativa, por ejemplo\
$ g(x) = -2x + 5$, da $ g(v) - g(u) = -2(v-u) < 0$: estrictamente \omterm{def:g10:functions:variations}{decreciente}.
\end{example}

\section{Extremos}

\begin{definition}[Máximo, mínimo]\label{def:g10:functions:extrema}
Sea $ f $ ser definido en $ D $ y $ a \in D $. el valor $ f(a)$ es el
\emph{máximo}\index{máximo} de $ f $ en $ D $ cuando $ f(x) \leq f(a)$ para todo
$ x \in D $; es el \emph{mínimo}\index{mínimo} cuando
$ f(x) \geq f(a)$ para todo $ x \in D $. en un grafo, son los más altos y
puntos más bajos de la curva.
\end{definition}

\begin{example}\label{ex:g10:functions:extremum}
mostrar que $ f(x) = (x - 1)^2 + 2$ tiene \omterm{def:g10:functions:extrema}{mínimo} $2$ en $\R $, alcanzado en
$ x = 1$. por cada $ x $, la plaza $(x-1)^2$ es $\geq 0$, entonces
$ f(x) = (x-1)^2 + 2 \geq 2$; y $ f(1) = 0 + 2 = 2$, entonces el valor $2$ es
realmente alcanzado. Ambos hechos juntos prueban que $2$ es el \omterm{def:g10:functions:extrema}{mínimo}.
\end{example}

\begin{omfigure}
\begin{tikzpicture}
\begin{axis}[omaxis,
  width=7.6cm, height=5.2cm,
  xmin=-1.8, xmax=3.8, ymin=-0.9, ymax=6.4,
  xtick={1}, ytick={2}]
  \addplot[omDef, thick, domain=-1.0:3.0, samples=100] {(x-1)^2 + 2};
  \addplot[omThm, dashed, domain=-1.6:3.6] {2};
  \fill (axis cs:1,2) circle (1.5pt);
  \node[below right, font=\small] at (axis cs:1,2) {$(1,2)$};
\end{axis}
\end{tikzpicture}

{\small la curva de $ f(x) = (x-1)^2 + 2$ nunca pasa por debajo de la línea discontinua
$ y = 2$, y lo toca en $ x = 1$: el \omterm{def:g10:functions:extrema}{mínimo} de $ f $ es $2$.}
\end{omfigure}

\section{Ejercicios}

\begin{exercise}[$\star $]\label{exo:g10:functions:1}
Sea $ f(x) = 2x^2 - 3x + 1$. Calcular el imágenes de $0$, $2$, $-1$ y
$\frac12$.
\end{exercise}

\begin{exercise}[$\star $]\label{exo:g10:functions:2}
Dale el \omterm{def:g10:functions:function}{dominio} de cada uno \omterm{def:g10:functions:function}{función}:
\[
f(x) = 5x - 2, \qquad
g(x) = \frac{3}{x + 4}, \qquad
h(x) = \sqrt{2x - 6}, \qquad
k(x) = \frac{1}{x^2 + 1}.
\]
\end{exercise}

\begin{exercise}[$\star $]\label{exo:g10:functions:3}
Sea $ f(x) = x^2 - 2x $. encontrar todo preimágenes de $0$, de $3$, y de $-1$.
\end{exercise}

\begin{exercise}[$\star $]\label{exo:g10:functions:4}
¿El punto $ A(-1, 4)$ pertenecer a la grafo de $ f(x) = x^2 - 2x + 1$?
y el punto $ B(2, 1)$? Justifíquelo mediante un cálculo.
\end{exercise}

\begin{exercise}[$\star $]\label{exo:g10:functions:5}
A \omterm{def:g10:functions:function}{función} $ g $ definido en $\intcc{-3}{5}$ tiene el \omterm{def:g10:functions:table}{tabla de variaciones}
\begin{center}
\renewcommand{\arraystretch}{1.3}
\begin{tabular}{c|ccccccc}
$ x $ & $-3$ & & $0$ & & $3$ & & $5$ \\
\hline
$ g $ & $1$ & $\searrow $ & $-2$ & $\nearrow $ & $4$ & $\searrow $ & $0$ \\
\end{tabular}
\end{center}
\begin{enumerate}
  \item ¿Cuáles son los \omterm{def:g10:functions:extrema}{máximo} y el \omterm{def:g10:functions:extrema}{mínimo} de $ g $ en $\intcc{-3}{5}$?
  \item Comparar $ g(-1)$ y $ g(-0.5)$ sin conocer sus valores.
  \item ¿Cuántas soluciones tiene el \omterm{def:g10:algebra:equation}{ecuación} $ g(x) = 0$ tener como \omterm{def:g10:functions:extrema}{máximo} en
        cada uno \omterm{def:g10:numbers:interval}{intervalo} de la mesa?
\end{enumerate}
\end{exercise}

\begin{exercise}[$\star\star $]\label{exo:g10:functions:6}
Sea $ f(x) = -2x + 7$. Mostrar, comparando $ f(u)$ y $ f(v)$ para $ u < v $,
eso $ f $ es estrictamente \omterm{def:g10:functions:variations}{decreciente} en $\R $.
\end{exercise}

\begin{exercise}[$\star\star $]\label{exo:g10:functions:7}
Demuestre que la \omterm{def:g10:functions:function}{función} $ f(x) = x^2 + 6x + 11$ tiene un \omterm{def:g10:functions:extrema}{mínimo} en $\R $, y
dar su valor y dónde se alcanza. (Pista: escribe
$ f(x) = (x + 3)^2 + c $ para la constante correcta $ c $.)
\end{exercise}

\begin{exercise}[$\star\star $]\label{exo:g10:functions:8}
Un jardín rectangular tiene perímetro. $40$ metro. Sea $ x $ ser su ancho, en
metros.
\begin{enumerate}
  \item Expresa la longitud y luego el área. $ A(x)$, como funciones de $ x $.
        para cual $ x $ ¿Tiene esto sentido?
  \item Calcular $ A(4)$, $ A(8)$, $ A(10)$, $ A(12)$ y $ A(16)$. ¿Qué haces?
        ¿Conjetura sobre la forma del área más grande?
\end{enumerate}
\end{exercise}

\begin{exercise}[$\star\star $]\label{exo:g10:functions:9}
Sea $ f(x) = \dfrac{1}{x^2 + 1}$, definido en $\R $.
\begin{enumerate}
  \item mostrar que $0 < f(x) \leq 1$ por cada real $ x $.
  \item Hace $ f $ tener un \omterm{def:g10:functions:extrema}{máximo} en $\R $? A \omterm{def:g10:functions:extrema}{mínimo}? Justificar.
\end{enumerate}
\end{exercise}

\begin{exercise}[$\star\star\star $]\label{exo:g10:functions:10}
Sea $ f(x) = x^2$ en $\R $.
\begin{enumerate}
  \item Sea $0 \leq u < v $. Factor $ f(v) - f(u)$ y deducir que $ f $ es
        estrictamente \omterm{def:g10:functions:variations}{creciente} en $\intco{0}{+\infty}$.
  \item Adapte el argumento para demostrar que $ f $ es estrictamente \omterm{def:g10:functions:variations}{decreciente} en
        $\intoc{-\infty}{0}$.
\end{enumerate}
\end{exercise}

\section{Problema: las máquinas alimentan su propia producción}

\begin{problem}[{Problema de fin de semana --- puntos fijos, iteración y
tres máquinas famosas: la mejoradora de Heron, la escuadradora y la
granizo sin resolver}]
\label{pb:g10:functions:1}
A \omterm{def:g10:functions:function}{función} Es una máquina: entra un número, sale un número.
Los experimentos más interesantes alimentan la máquina. \emph{su propio
salida}, una y otra vez --- y luego dominan dos preguntas:
¿El proceso se asienta en algún lugar (un \emph{\omterm{pb:g10:functions:1}{punto fijo}}), y
¿Cómo la lleva la variación de la máquina hasta allí? uno de esto
La máquina del problema ha estado puliendo raíces cuadradas durante dos
mil años; otro guarda el problema sin resolver más famoso
las matemáticas elementales tienen para ofrecer.

\medskip
\noindent\textbf{Part I --- Heron's machine.}
Considere la \omterm{def:g10:functions:function}{función}
\[
f(x) = \frac12\left(x + \frac2x\right).
\]
\begin{enumerate}
  \item Dale el \omterm{def:g10:functions:function}{dominio} de $ f $y calcular $ f(1)$, $ f(2)$ y
        $ f\!\left(\frac32\right)$ como fracciones exactas.
  \item encontrar todo preimágenes de $\frac32$: resolver
        $ f(x) = \frac32$ (borra el denominador y usa el
        métodos de \cref{pb:g10:algebra:1}).
  \item A \emph{punto fijo}\index{punto fijo} de $ f $ es un
        número no modificado por la máquina: $ f(x) = x $. encontrar ambos
        puntos fijos de $ f $. ¿Qué viejo conocido guarda el
        uno positivo?
  \item A partir de $ x = 1$, alimenta la máquina con su propia salida
        tres veces y registrar los resultados exactos. ¿Qué sucesión
        del problema de la irracionalidad del fin de semana del volumen de la Escuela Secundaria, ¿has reconstruido --- y lo que fue
        ``La receta de la garza'', en el lenguaje de funciones?
  \item En una \omterm{def:g10:functions:function}{imagen}, dibuja el grafo de $ f $ (para $ x > 0$)
        y la linea diagonal $ y = x $. donde estan los fijo
        puntos en esta foto? (Comparar
        El problema del fin de semana de dos termómetros del volumen de la escuela secundaria: la lectura de temperatura es la misma
        en ambas escalas fue la misma idea.)
\end{enumerate}

\medskip
\noindent\textbf{Part II --- Why the machine cannot miss.}
\begin{enumerate}[resume]
  \item Para $0 < u < v $, mostrar por reducción a un común
        denominador que
        \[
        f(v) - f(u) = \frac{(v - u)(uv - 2)}{2uv} ,
        \]
        y deducir las variaciones de $ f $ en
        $\intoo{0}{+\infty}$: \omterm{def:g10:functions:variations}{decreciente} hasta cierto punto entonces
        \omterm{def:g10:functions:variations}{creciente} --- ¿Qué punto? (La misma técnica que
        \cref{exo:g10:functions:10}.)
  \item Deduce eso en $\intoo{0}{+\infty}$ la \omterm{def:g10:functions:function}{función} $ f $ tiene
        un \omterm{def:g10:functions:extrema}{mínimo}, vale exactamente $\sqrt2$, alcanzado en
        $\sqrt2$. ¿Qué dice esto sobre \emph{cada} salida
        de la máquina (entradas positivas)?
  \item Demuestre que si $ x > \sqrt2$ entonces
        $\sqrt2 < f(x) < x $: una suposición demasiado grande siempre mejora,
        nunca se sobrepasó. (Para $ f(x) < x $, comparar $\frac2x $ con
        $ x $.)
  \item Reúna las preguntas 6 a 8 en el \omterm{def:g10:functions:table}{tabla de variaciones} de $ f $
        en $\intoo{0}{+\infty}$
        (\cref{def:g10:functions:table}), \omterm{def:g10:functions:extrema}{mínimo} incluido.
  \item Una segunda máquina: $ g(x) = x^2$. Es puntos fijos son
        $0$ y $1$ (controlar). Iterar $ g $ cuatro veces desde
        $ x = 0.9$, luego cuatro veces desde $ x = 1.1$ (calculadora,
        tres decimales). Uno \omterm{pb:g10:functions:1}{punto fijo} atrae, el otro
        repele: ¿cuál es cuál?
\end{enumerate}

\medskip
\noindent\textbf{Part III --- The hailstone machine.}
Sobre los números enteros, defina: $ h(n) = \frac n2$ si $ n $ es par,
y $ h(n) = 3n + 1$ si $ n $ es impar. Los números rebotan $ h $ me gusta
Granizo en una nube de tormenta.
\begin{enumerate}[resume]
  \item Calcular el vuelo completo de $7$ bajo $ h $, hasta $1$.
        ¿Cuántos pasos da y a qué altura vuela?
  \item Iniciar el vuelo de $27$ y calcular diez pasos. (Su
        dura el vuelo completo $111$ pasos y picos en $9\,232$ ---
        desde un punto de partida de $27$.) ¿Qué lección sobre simple?
        las reglas lo hacen $7$ y $27$ ¿enseñar?
  \item Explique por qué cualquier vuelo que alcance una potencia de $2$
        se estrella directamente a $1$, y qué sucede en $1$
        (calcular $ h(1)$, $ h(4)$, $ h(2)$). Hace $ h $ tener alguna
        \omterm{pb:g10:functions:1}{punto fijo} entre los números enteros positivos? ¿Qué hace?
        ¿El final de cada vuelo observado?
  \item El \emph{Conjetura de Collatz} reclama cada inicio
        el número finalmente llega $1$. Las computadoras han verificado
        mucho más allá $10^{20}$; ningún humano lo ha demostrado. que
        dos tipos de descubrimiento \emph{refutar} ¿él? y por que
        ¿La verificación masiva aún no cierra el caso?
        (¿El problema del fin de semana del oráculo de cerillas del volumen de la escuela secundaria hizo sonar esta campana)?
  \item "Las matemáticas aún no están maduras para tales cuestiones", dijo
        Pablo Erd\H{o}s de esta conjetura. En una frase: ¿qué?
        separa la máquina de granizo de la de Heron, donde
        ¿Las preguntas 6 a 8 resolvieron todo?
\end{enumerate}

\medskip
\noindent\textbf{Part IV --- The machine user's manual.}
\begin{enumerate}[resume]
  \item \omterm{def:g10:functions:function}{Dominios} como \omterm{def:g10:numbers:interval}{intervalos} (\cref{pb:g10:numbers:1}'s
        idioma): dar el \omterm{def:g10:functions:function}{dominios} de
        $ x \mapsto \sqrt{x - 3}$, \
        $ x \mapsto \frac{1}{x^2 - 9}$, \
        $ x \mapsto \sqrt{9 - x^2}$.
  \item Volver a Heron $ f $: encontrar el exacto preimágenes de $3$
        (resolver $ f(x) = 3$ completando el cuadrado).
  \item Se lanza una piedra hacia arriba: su altura es
        $ H(t) = 20t - 5t^2$ metros después $ t $ artículos de segunda clase. en cual
        tiempo \omterm{def:g10:numbers:interval}{intervalo} ¿Tiene sentido físico la fórmula?
        Completa el cuadrado para encontrar la altura máxima y su
        tiempo (\cref{def:g10:functions:extrema}), y dar la
        \omterm{def:g10:functions:table}{tabla de variaciones}.
  \item ¿Cuánto tiempo dura al menos la piedra? $15$ ¿m alto? (Resolver
        $ H(t) \geq 15$ con un \omterm{def:g10:algebra:signtable}{tabla de signos},
        \cref{met:g10:algebra:signtable}.)
  \item Final: este problema ejecutó cuatro máquinas --- Heron's
        mejorador, el escuadrador, el granizo, la piedra
        altura. En dos o tres oraciones, indica lo que cada
        \omterm{def:g10:functions:function}{función} tiene (un \omterm{def:g10:functions:function}{dominio}, una regla, una grafo) y los dos
        grandes preguntas que este problema planteó de cada uno (¿dónde
        ¿Se resuelve la iteración? ¿Cómo funciona la \omterm{def:g10:functions:function}{función} variar?) ---
        preguntas que los próximos capítulos abordan una por una.
\end{enumerate}
\end{problem}
